\documentclass[11pt,letterpaper]{article}

\usepackage[spanish,es-noshorthands]{babel}
\usepackage[T1]{fontenc}
\usepackage[utf8]{inputenc}
\usepackage{lmodern}
\usepackage{microtype}
\usepackage{geometry}
\usepackage{booktabs}
\usepackage{tabularx}
\usepackage{enumitem}
\usepackage{xcolor}
\usepackage{csquotes}
\usepackage[hidelinks]{hyperref}
\usepackage[backend=biber,style=apa,sorting=nyt]{biblatex}
\DeclareLanguageMapping{spanish}{spanish-apa}

\geometry{margin=2.5cm}
\setlength{\parindent}{0pt}
\setlength{\parskip}{0.65em}
\setlist{nosep}
\addbibresource{referencias.bib}

\newcommand{\pendiente}[1]{\textcolor{gray}{[#1]}}

\title{\textbf{Innovaciones inspiradas en la naturaleza}\\
\large CRISPR-Cas9: de la inmunidad bacteriana a la edición genética}
\author{\pendiente{Nombre del estudiante}\\
Gestión del Desarrollo Tecnológico}
\date{\pendiente{Fecha de entrega}}

\begin{document}

\maketitle

\section{Introducción}

La naturaleza puede entenderse como un extenso laboratorio de investigación y
desarrollo. A través de la evolución, los organismos han generado mecanismos
para obtener energía, defenderse, comunicarse y adaptarse a condiciones
cambiantes. Cuando un diseño humano estudia esos mecanismos y traslada sus
principios a una solución técnica, se produce una innovación bioinspirada. La
transferencia no tiene que limitarse a copiar la forma visible de un organismo;
también puede aprovechar procesos que ocurren en las escalas celular y
molecular.

CRISPR-Cas9 constituye un caso especialmente profundo de esta transferencia.
La tecnología no surgió inicialmente como una herramienta construida para
editar el genoma humano. Su antecedente es un sistema mediante el cual ciertas
bacterias y arqueas reconocen material genético invasor y se protegen de
nuevas infecciones. La investigación científica permitió comprender ese
proceso natural, simplificarlo y reprogramarlo para cortar ADN en lugares
elegidos por una persona. Este caso muestra cómo el conocimiento de una función
biológica puede convertirse en una plataforma tecnológica con aplicaciones en
investigación, medicina y agricultura.

\section{CRISPR-Cas9: una herramienta inspirada en la inmunidad bacteriana}

\subsection{Antecedente en la naturaleza}

Las bacterias están expuestas a bacteriófagos, virus que introducen su material
genético dentro de ellas para reproducirse. Algunas bacterias responden a esta
amenaza mediante el sistema CRISPR-Cas, una forma de inmunidad adaptativa. El
acrónimo CRISPR designa unas secuencias repetidas del ADN bacteriano separadas
por fragmentos variables llamados \emph{espaciadores}. Muchos de estos
espaciadores proceden del ADN de virus encontrados anteriormente.

El proceso natural puede resumirse en tres funciones. Primero, durante la
\textbf{adaptación}, la bacteria incorpora un fragmento del ADN invasor en su
región CRISPR y conserva así un registro molecular de la infección. Después,
durante la \textbf{expresión}, esa región se transcribe y procesa para producir
pequeños ARN CRISPR (crRNA). Finalmente, durante la
\textbf{interferencia}, el crRNA guía a proteínas asociadas con CRISPR hacia
una secuencia complementaria del material invasor, que es reconocida y
cortada. En 2007, Barrangou y sus colaboradores demostraron
experimentalmente que la adquisición de nuevos espaciadores podía conferir
resistencia frente a bacteriófagos y que la especificidad dependía de la
correspondencia entre el espaciador y la secuencia viral
\parencite{barrangou2007}.

El sistema tipo II de la bacteria \emph{Streptococcus pyogenes} utiliza la
proteína Cas9. En él participa además un ARN transactivador, o tracrRNA, que
contribuye a la maduración del crRNA y es esencial para la defensa contra ADN
invasor \parencite{deltcheva2011}. El crRNA aporta la información para
reconocer el objetivo, mientras Cas9 actúa como una endonucleasa: una enzima
capaz de cortar las dos hebras del ADN. La bacteria combina, por tanto,
\emph{memoria}, \emph{reconocimiento} y \emph{respuesta} en un sistema
molecular programable.

\subsection{De fenómeno natural a diseño tecnológico}

El paso decisivo consistió en abstraer el principio esencial del sistema:
una secuencia de ARN puede dirigir una enzima de corte hacia una secuencia
específica de ADN. En 2011, el equipo de Emmanuelle Charpentier describió la
función del tracrRNA \parencite{deltcheva2011}. Posteriormente, Charpentier y
Jennifer Doudna colaboraron para reconstruir y simplificar el mecanismo.

En el trabajo publicado por Jinek y sus colaboradores en 2012, se demostró que
Cas9 podía programarse para cortar secuencias seleccionadas de ADN. Además, el
crRNA y el tracrRNA naturales se fusionaron en una sola molécula sintética,
denominada ARN guía único. Al modificar la parte variable de esta guía, el
investigador puede dirigir Cas9 hacia un nuevo objetivo sin tener que diseñar
una proteína distinta para cada secuencia \parencite{jinek2012}.

La analogía de las ``tijeras genéticas'' resulta útil, pero es incompleta. El
sistema no reescribe por sí solo el ADN: Cas9 genera un corte y la célula lo
repara. Una reparación directa puede introducir pequeñas inserciones o
eliminaciones y desactivar un gen. Si se aporta una plantilla adecuada, otros
mecanismos celulares pueden incorporar un cambio específico. Por consiguiente,
la edición depende tanto de la precisión del reconocimiento como de la vía de
reparación que se active.

\begin{table}[ht]
  \centering
  \small
  \caption{Traducción del sistema natural a la herramienta tecnológica.}
  \label{tab:traduccion-crispr}
  \begin{tabularx}{\textwidth}{@{}>{\raggedright\arraybackslash}X
    >{\raggedright\arraybackslash}X@{}}
    \toprule
    \textbf{Sistema natural} & \textbf{Diseño tecnológico} \\
    \midrule
    El espaciador conserva información de un virus previo.
      & La persona diseña una secuencia guía para el objetivo elegido. \\
    El crRNA y el tracrRNA participan en la localización del ADN invasor.
      & Ambos componentes se simplifican en un ARN guía único. \\
    Cas9 corta ADN viral complementario y protege a la bacteria.
      & Cas9 corta una región seleccionada del genoma para permitir su edición. \\
    La finalidad biológica es la supervivencia frente a una infección.
      & La finalidad humana es estudiar, desactivar, corregir o modificar genes. \\
    \bottomrule
  \end{tabularx}
\end{table}

\subsection{Objetivos y propuesta de valor}

El objetivo inicial del diseño de 2012 fue demostrar una endonucleasa de ADN
programable mediante ARN, no presentar de inmediato una terapia terminada.
Como plataforma de investigación, su propuesta de valor consistía en separar
dos funciones: Cas9 conservaba la capacidad de corte y el ARN guía determinaba
el destino. Esta modularidad hizo que cambiar de objetivo fuera relativamente
rápido y redujo la complejidad de diseñar herramientas específicas.

Las finalidades posteriores pueden agruparse en tres ámbitos:

\begin{itemize}
  \item \textbf{Investigación básica:} desactivar o modificar genes para
    estudiar su función y construir modelos celulares o animales.
  \item \textbf{Salud:} explorar tratamientos para enfermedades de base
    genética, modificar células con fines terapéuticos y desarrollar nuevas
    estrategias contra el cáncer.
  \item \textbf{Agricultura y biotecnología:} generar cultivos con propiedades
    seleccionadas y modificar microorganismos utilizados en procesos
    productivos.
\end{itemize}

\subsection{Proceso y contexto de lanzamiento}

CRISPR-Cas9 no tuvo un lanzamiento convencional como producto de consumo. Su
introducción ocurrió mediante publicaciones científicas que convirtieron un
hallazgo biológico en una plataforma reproducible. La trayectoria comenzó con
la observación, en 1987, de una estructura de secuencias repetidas cuyo
significado todavía era desconocido \parencite{ishino1987}. Décadas de
investigación posterior relacionaron los espaciadores con virus y demostraron
la función inmunitaria del sistema \parencite{barrangou2007}.

El artículo de 2012 representó el punto de inflexión tecnológico: mostró que
Cas9 podía recibir instrucciones sintéticas para cortar ADN en sitios
predeterminados \parencite{jinek2012}. En 2013, distintos equipos demostraron
su aplicación para editar células de mamíferos y humanas
\parencite{cong2013,mali2013}. Esta rápida adopción indica que el lanzamiento
efectivo se produjo dentro de una comunidad científica capaz de reproducir,
mejorar y ampliar el diseño.

En 2020, la Real Academia Sueca de Ciencias otorgó a Charpentier y Doudna el
Premio Nobel de Química ``por el desarrollo de un método para la edición del
genoma'' \parencite{nobel2020}. El reconocimiento no correspondió al
descubrimiento inicial de las secuencias CRISPR, sino a la transformación del
mecanismo en una herramienta programable. Esta distinción evidencia el
carácter acumulativo de la innovación: observaciones, explicación funcional,
simplificación de componentes y adopción tecnológica fueron etapas
complementarias.

\subsection{Impacto, límites y gobernanza}

La misma capacidad de programación que hace valiosa a CRISPR-Cas9 genera
desafíos. El ARN guía puede conducir cortes en sitios no deseados; la reparación
celular no siempre produce el resultado previsto; y llevar los componentes
hasta el tejido correcto continúa siendo un problema técnico. Además, las
consecuencias éticas varían según se editen células somáticas de una persona o
células reproductivas cuyos cambios podrían heredarse.

Por ello, el desarrollo tecnológico no puede evaluarse únicamente por su
eficacia molecular. También requiere gobernanza, supervisión, acceso
equitativo y deliberación sobre usos aceptables. La Organización Mundial de la
Salud ha propuesto recomendaciones para la supervisión de la edición del
genoma humano ante sus implicaciones científicas, éticas, sociales y legales
\parencite{who2021}. Desde la gestión tecnológica, CRISPR-Cas9 demuestra que
una plataforma de gran alcance necesita evolucionar junto con las capacidades
institucionales que controlan su aplicación.

\subsection{Síntesis del caso}

CRISPR-Cas9 es un diseño inspirado en la biología porque transfiere un
principio funcional de la inmunidad bacteriana a un contexto humano. La
naturaleza aportó un sistema que registra amenazas, reconoce secuencias y
corta ADN invasor. La ingeniería conservó la lógica de reconocimiento y corte,
pero sustituyó el objetivo evolutivo de defensa por objetivos definidos por el
investigador. La innovación central no consistió en copiar la forma de un
organismo, sino en comprender, simplificar y reprogramar un proceso molecular.

\section{Algoritmos de optimización inspirados en el moho mucilaginoso}

\subsection{Antecedente en la naturaleza}

\emph{Physarum polycephalum} es un moho mucilaginoso plasmodial, perteneciente
al grupo de los amebozoos y no al de los hongos verdaderos. Durante una etapa
de su ciclo de vida forma una gran célula multinucleada que se desplaza sobre
superficies húmedas en busca de alimento. Su cuerpo adopta la forma de una red
de tubos por los que circula el protoplasma. Esta red transporta nutrientes y
señales entre las distintas regiones del organismo.

Aunque \emph{Physarum} no posee cerebro ni sistema nervioso, su estructura
cambia de acuerdo con las condiciones del entorno. Cuando encuentra varias
fuentes de alimento, inicialmente extiende numerosas conexiones. Con el tiempo
refuerza algunos tubos y elimina otros, manteniendo rutas que permiten
transportar nutrientes con un costo material reducido. No se trata de que el
organismo construya o compare conscientemente un conjunto de mapas. La
solución emerge de interacciones locales entre flujo, contracciones y grosor
de los tubos.

Una demostración influyente fue publicada por Nakagaki, Yamada y Tóth en 2000.
Los investigadores distribuyeron el plasmodio dentro de un laberinto y
colocaron alimento en dos puntos. El organismo ocupó inicialmente varios
caminos, pero después conservó una conexión cercana a la ruta de menor longitud
entre las dos fuentes \parencite{nakagaki2000}. El experimento mostró que un
sistema biológico sin control central podía resolver físicamente un problema
de optimización.

\subsection{Principio biológico esencial}

El principio transferible es un ciclo de \textbf{exploración, realimentación y
poda}. Durante la exploración se forma una red con rutas alternativas. El flujo
de protoplasma no se distribuye por igual: depende de la longitud y la
conductividad de cada tubo. Una ruta por la que circula más flujo recibe
realimentación positiva y aumenta su capacidad. Las rutas con poco flujo se
debilitan progresivamente hasta desaparecer.

Esta regla local produce una propiedad global. La red tiende a reducir el
material necesario para conectar las fuentes de alimento, pero puede mantener
conexiones redundantes que mejoran la tolerancia a interrupciones. Por ello, el
resultado no debe interpretarse solamente como el camino geométricamente más
corto. En redes con múltiples destinos, el organismo negocia varios criterios:
longitud total, eficiencia del transporte y resiliencia.

El diseño es relevante para la gestión tecnológica porque utiliza
\textbf{control descentralizado}. Ningún componente individual conoce la
solución completa; la configuración se modifica a partir de información local.
Además, la red continúa adaptándose cuando aparece alimento, desaparece una
fuente o cambia el ambiente. Estos atributos resultan atractivos para sistemas
humanos que operan bajo incertidumbre y no pueden depender siempre de un
controlador central.

\subsection{Del organismo al modelo matemático}

Tero, Kobayashi y Nakagaki formalizaron en 2007 un modelo de red de transporte
adaptativa basado en las observaciones fisiológicas de \emph{Physarum}
\parencite{tero2007}. En su traducción computacional, el problema se representa
mediante un grafo: las fuentes de alimento se convierten en nodos y los
posibles tubos en aristas. Cada arista posee una longitud y una conductividad.
Una diferencia de presión produce un flujo simulado a través de la red.

El algoritmo repite dos operaciones centrales:

\begin{enumerate}
  \item calcula cómo se distribuiría el flujo entre las conexiones existentes;
  \item aumenta la conductividad de las conexiones con mayor flujo y reduce la
    de aquellas que reciben poco.
\end{enumerate}

Después de varias iteraciones, las rutas ineficientes pierden conductividad y
las más útiles se estabilizan. El modelo convierte así la morfología dinámica
del organismo en una regla de actualización numérica. Su importancia no reside
en simular cada detalle celular, sino en conservar la relación funcional:
\emph{mayor uso produce refuerzo; menor uso produce desaparición}.

\begin{table}[ht]
  \centering
  \small
  \caption{Traducción del comportamiento de \emph{Physarum} a un modelo de
  optimización.}
  \label{tab:traduccion-physarum}
  \begin{tabularx}{\textwidth}{@{}>{\raggedright\arraybackslash}X
    >{\raggedright\arraybackslash}X@{}}
    \toprule
    \textbf{Sistema natural} & \textbf{Modelo computacional} \\
    \midrule
    Las fuentes de alimento atraen al organismo.
      & Los destinos se representan como nodos de un grafo. \\
    Los tubos conectan distintas regiones del plasmodio.
      & Las conexiones posibles se representan mediante aristas. \\
    El protoplasma circula por los tubos.
      & Se calcula un flujo a partir de presión, longitud y conductividad. \\
    Los tubos con mayor flujo se refuerzan.
      & La conductividad de las aristas útiles aumenta. \\
    Los tubos poco utilizados se atrofian.
      & Las conexiones ineficientes pierden peso o se eliminan. \\
    La red se reorganiza ante cambios ambientales.
      & El algoritmo recalcula la solución cuando cambian los datos. \\
    \bottomrule
  \end{tabularx}
\end{table}

\subsection{Objetivos y aplicaciones}

El objetivo tecnológico es encontrar configuraciones eficientes en problemas
con numerosas soluciones posibles. En su forma más directa, el modelo permite
buscar caminos cortos dentro de un grafo. Sin embargo, su valor aumenta cuando
el problema exige conectar varios puntos y equilibrar criterios que pueden
competir entre sí.

Entre sus contextos de aplicación se encuentran:

\begin{itemize}
  \item diseño de redes de transporte, carreteras y ferrocarriles;
  \item enrutamiento de información en redes de comunicación;
  \item planificación de cadenas de suministro y redes logísticas;
  \item aproximación de problemas de conexión y caminos mínimos;
  \item optimización de parámetros en problemas de ingeniería.
\end{itemize}

Estas aplicaciones no implican introducir organismos vivos en la
infraestructura. El producto transferido es la regla de adaptación obtenida al
estudiar su comportamiento. En otros experimentos, el organismo físico también
puede utilizarse como una forma de computación no convencional, pero el diseño
que aquí se analiza es principalmente matemático y computacional.

\subsection{Contexto de desarrollo y lanzamiento}

Este caso tampoco tuvo un único lanzamiento comercial. Su evolución ocurrió a
través de una secuencia de publicaciones que trasladaron el conocimiento desde
la biología experimental hasta la optimización.

El experimento del laberinto publicado en 2000 aportó una evidencia visible de
selección de rutas \parencite{nakagaki2000}. El modelo matemático de 2007
identificó el mecanismo de realimentación entre flujo y grosor del tubo y lo
convirtió en una dinámica aplicable a grafos \parencite{tero2007}. En 2010,
Tero y sus colaboradores ampliaron el principio al diseño de redes adaptativas.
Para evaluar el modelo, colocaron fuentes de alimento en posiciones
correspondientes a ciudades del área de Tokio. Las redes producidas por
\emph{Physarum} y por el modelo alcanzaron combinaciones de costo, eficiencia y
tolerancia a fallos comparables con características de la red ferroviaria real
\parencite{tero2010}.

El contexto de lanzamiento fue, por tanto, interdisciplinario. Participaron la
biología, la física, las matemáticas y la planificación de redes. La innovación
no buscó reemplazar inmediatamente a los algoritmos clásicos, sino demostrar
que una regla biológica sencilla podía generar diseños competitivos y
adaptativos. Su propuesta de valor se encontraba en la combinación de
descentralización, flexibilidad y robustez.

\subsection{El \emph{Slime Mould Algorithm} de 2020}

Una línea posterior de desarrollo es el \emph{Slime Mould Algorithm} (SMA),
presentado por Li y colaboradores en 2020 como un método de optimización
estocástica \parencite{li2020}. Aunque comparte la inspiración biológica, no
debe confundirse con el modelo de flujo de Tero. El SMA pertenece a la familia
de metaheurísticas poblacionales: mantiene múltiples soluciones candidatas y
actualiza sus posiciones dentro de un espacio de búsqueda.

El algoritmo abstrae tres aspectos del forrajeo de \emph{Physarum}: la
aproximación hacia zonas con mayor concentración de alimento, las oscilaciones
del organismo y la realimentación positiva o negativa según la calidad de la
fuente encontrada. Las soluciones con mejor desempeño reciben pesos mayores y
orientan la búsqueda, mientras componentes aleatorios permiten explorar otras
regiones y reducir el riesgo de quedar atrapado prematuramente en una solución
local. El trabajo de lanzamiento evaluó el método sobre funciones de prueba y
problemas de ingeniería con restricciones.

La distinción entre ambos desarrollos fortalece el análisis del caso. El modelo
de redes adaptativas reproduce el refuerzo de conductos y es especialmente
intuitivo para rutas y transporte. El SMA toma una inspiración más general del
forrajeo y la transforma en una estrategia de optimización aplicable a
variables continuas. Ambos proceden del mismo antecedente natural, pero
responden a objetivos y arquitecturas algorítmicas diferentes.

\subsection{Impacto, límites y evaluación crítica}

La principal contribución de \emph{Physarum} al diseño tecnológico es mostrar
que la eficiencia de una red puede emerger de reglas locales y adaptación
continua. Esto abre alternativas a los modelos que requieren conocer de
antemano toda la estructura del problema. También ofrece una manera de explorar
soluciones que equilibran costo y resiliencia, en lugar de optimizar una sola
métrica.

No obstante, la analogía biológica no garantiza superioridad computacional.
Los modelos simplifican el comportamiento del organismo y su desempeño depende
de parámetros, condiciones iniciales y criterios de comparación. Las
metaheurísticas, incluido el SMA, pueden encontrar soluciones de alta calidad,
pero no necesariamente demuestran que hayan alcanzado el óptimo global. Además,
los algoritmos clásicos pueden resultar más rápidos o previsibles en problemas
bien definidos.

Por ello, la evaluación debe comparar resultados, tiempo de cómputo,
estabilidad y facilidad de implementación, no limitarse al atractivo de la
metáfora natural. La bioinspiración aporta una hipótesis de diseño; la
validación experimental y computacional determina si esa hipótesis crea valor
en un contexto específico.

\subsection{Síntesis del caso}

Los algoritmos inspirados en \emph{Physarum polycephalum} trasladan a la
computación una estrategia natural de formación de redes. El organismo explora
el ambiente, distribuye flujo, refuerza conexiones productivas y elimina rutas
poco utilizadas. Los diseños humanos convierten esas acciones en nodos,
aristas, conductividades, pesos y reglas iterativas. El resultado es una
familia de métodos que busca soluciones eficientes sin depender de un plan
centralizado.

En comparación con CRISPR-Cas9, la transferencia ocurre en otra escala.
CRISPR adapta directamente componentes moleculares de un sistema inmunitario;
los algoritmos de \emph{Physarum} abstraen un comportamiento y lo expresan
mediante matemáticas. Los dos casos muestran que la innovación bioinspirada
puede surgir tanto de reutilizar mecanismos materiales de la naturaleza como
de convertir sus principios organizativos en información ejecutable.

% ---------------------------------------------------------------------------
% ESTRUCTURA PARA LAS SIGUIENTES ITERACIONES
% Descomentar después de seleccionar los otros casos.
%
% \section{Tercer diseño bioinspirado}
% \subsection{Antecedente en la naturaleza}
% \subsection{Principio esencial y traducción tecnológica}
% \subsection{Objetivos, proceso y contexto de lanzamiento}
% \subsection{Impacto y limitaciones}
%
% \section{Comparación de los casos}
% Comparar tipo de inspiración, problema abordado, proceso de transferencia,
% contexto de lanzamiento, impacto y límites.
%
% \section{Conclusiones}
% Integrar los hallazgos de los dos o tres casos.
% ---------------------------------------------------------------------------

\printbibliography[title={Referencias}]

\end{document}
